% GENERATED FILE. Edit canonical lessons, then run npm run generate:books.
% canonical-source-hash: fnv1a64:26a5eb14bdbc3854
% canonical-lessons: ZH-C313-zhuangxiu, ZH-C313-cai, ZH-W313-cai-only, ZH-C313-dingju, ZH-C313-shebude, ZH-C313-shiying

\chapter{Moving House}
\label{ch:zh-moving-house}
\hlchaptermodality{313}

\begin{chapteropening}
\textbf{By the end of this chapter:} \emph{I can talk about doing up a new home, settling down and getting used to a new place.}

Complete the last lesson of chapter 313: I can talk about doing up a new home, settling down and getting used to a new place.
\end{chapteropening}

\section[zhuāngxiū]{zhuāngxiū — to do up a home; to renovate}

\label{lesson:ZH-C313-zhuangxiu}



\begin{quote}
Before the new one: say the Mandarin for it serves them right, then the Mandarin for whatever anyone says.
\end{quote}

\subsection*{You'll want to know: zhuāngxiū}
\emph{zhuāngxiū} — \textquotedblleft{}to do up a home; to renovate\textquotedblright{}.

\begin{quote}
\emph{wǒmen de xīn jiā zhèngzài zhuāngxiū} — Our new home is being done up.
\end{quote}

Its characters come later; for now it is a word for the ear.

\begin{cousinweb}
\emph{zhuāngxiū} is built from \emph{zhuāng} 'to pack' and \emph{xiū} 'to repair', both inherited from Middle Chinese, \emph{tsrjang} and \emph{sjuw} (Guangyun, in Baxter's transcription), which glosses the first 'to pack for a journey' and the second 'to put in order'.
\end{cousinweb}

\begin{grammarlens}[title={What you've built}]
One more word for telling what happened.
\end{grammarlens}

\subsection*{Your turn}
\begin{itemize}
\raggedright
  \item \emph{Say it:} \emph{zhuāngxiū}
  \item \emph{Say it:} \emph{zhuāngxiū}, once more
  \item \emph{Say it:} \emph{wǒmen de xīn jiā zhèngzài zhuāngxiū} — Our new home is being done up.
  \item \emph{Recall:} say the Mandarin for it serves them right, then the Mandarin for whatever anyone says, then say \emph{zhuāngxiū} again
\end{itemize}

\begin{tcolorbox}[breakable,colback=teal!4,colframe=teal!35!black,title={Before you move on}]
What does zhuāngxiū mean? (\textquotedblleft{}To do up a home; to renovate\textquotedblright{}.) Say it once more.
\end{tcolorbox}
\section[\texorpdfstring{\zh{才} (cái)}{才 (cái)}]{\zh{才} (cái) — only then; not until}

\label{lesson:ZH-C313-cai}



\begin{quote}
Before the new one: say the Mandarin for whatever anyone says, then the Mandarin for to renovate a flat.
\end{quote}

\subsection*{You'll want to know: \zh{才}}
\textbf{\zh{才}} — \emph{cái} — \textquotedblleft{}only then; not until\textquotedblright{}.

\begin{quote}
\textbf{\zh{他十点才来}} \emph{tā shí diǎn cái lái} — He didn't come until ten.
\end{quote}

The new character, in pieces: \zh{才} — three strokes: a bar, a vertical that hooks to the left at the foot, and a long sweep down to the left across them: the hand piece of \zh{打} with a sweep for its last stroke. The next lesson puts it in your hand.

\begin{cousinweb}
\textbf{\zh{才}} (\emph{cái}) is inherited from Middle Chinese \emph{dzoj} (Guangyun, in Baxter's transcription), which glosses it 'barely' under another, more complex spelling; the simple form, read the same, is a different word there, glossed 'use', 'strength' and 'talent'.
\end{cousinweb}

\begin{grammarlens}[title={What you've built}]
One more word for telling what happened.
\end{grammarlens}

\subsection*{Your turn}
\begin{itemize}
\raggedright
  \item \emph{Say it:} \emph{cái}
  \item \emph{Say it:} \emph{cái}, once more
  \item \emph{Say it:} \emph{tā shí diǎn cái lái} — He didn't come until ten.
  \item \emph{Recall:} say the Mandarin for whatever anyone says, then the Mandarin for to renovate a flat, then say \emph{cái} again
\end{itemize}

\begin{tcolorbox}[breakable,colback=teal!4,colframe=teal!35!black,title={Before you move on}]
What does \zh{才} mean? (\textquotedblleft{}Only then; not until\textquotedblright{}.) Say it once more.
\end{tcolorbox}
\section[\texorpdfstring{\zh{才} (cái)}{才 (cái)}]{\zh{才} — a bar, a hooked post and a sweep}

\label{lesson:ZH-W313-cai-only}



\begin{quote}
The word first, then the character inside it.

\begin{itemize}
\raggedright
  \item \emph{Recall:} say \emph{cái} — \textbf{R1}, one lesson back, before it decays
\end{itemize}

\begin{quote}
\zh{打} \zh{十}
\end{quote}

\emph{Read it:} each one aloud, with nothing said first — \emph{dǎ}, \emph{shí}
\end{quote}

\subsection*{Script — a bar, a hooked post and a sweep}
\hlblockfigure{\detokenize{figures/ZH-W313-cai-only-filmstrip.pdf}}{How \zh{才} is written, stroke by stroke}

\begin{quote}
\zh{才}
\end{quote}

\textbf{One piece, three strokes.}

\textbf{A bar} across the top.

\textbf{A vertical} down through the bar, hooking to the left at the foot.

\textbf{A long sweep} from the crossing, down to the lower left, across the vertical. In \zh{打} this stroke rises to the right; here it falls to the left.

\subsection*{Your turn}
\begin{itemize}
\raggedright
  \item \emph{Write it:} \zh{才}, all three, with the model in front of you
  \item \emph{Cover:} write it again from memory, and count the pieces, not the strokes
\end{itemize}

\begin{tcolorbox}[breakable,colback=teal!4,colframe=teal!35!black,title={Before you move on}]
How many strokes in \zh{才}? (\textbf{Three}.) How is it different from the hand piece of \zh{打}? (\textbf{The last stroke sweeps down to the left}.)

Source: \href{https://raw.githubusercontent.com/chanind/hanzi-writer-data/68d10a4b21150cae5e1ebbd223eed289cf32d90c/data/\%E6\%89\%8D.json}{Hanzi Writer Data \zh{才}, pinned stroke paths and medians}.
\end{tcolorbox}
\section[dìngjū]{dìngjū — to settle down; to make a place your home}

\label{lesson:ZH-C313-dingju}



\begin{quote}
Before the new one: say the Mandarin for to renovate a flat, then the Mandarin for not until then.
\end{quote}

\subsection*{You'll want to know: dìngjū}
\emph{dìngjū} — \textquotedblleft{}to settle down; to make a place your home\textquotedblright{}.

\begin{quote}
\emph{tāmen zài Běijīng dìngjū le} — They've settled in Beijing.
\end{quote}

Its characters come later; for now it is a word for the ear.

\begin{cousinweb}
\emph{dìngjū} is built from \emph{dìng} 'calm' and \emph{jū} 'to dwell', both inherited from Middle Chinese, \emph{dengH} and \emph{kjo} (Guangyun, in Baxter's transcription), which glosses the first 'at peace' and the second 'to dwell'.
\end{cousinweb}

\begin{grammarlens}[title={What you've built}]
One more word for telling what happened.
\end{grammarlens}

\subsection*{Your turn}
\begin{itemize}
\raggedright
  \item \emph{Say it:} \emph{dìngjū}
  \item \emph{Say it:} \emph{dìngjū}, once more
  \item \emph{Say it:} \emph{tāmen zài Běijīng dìngjū le} — They've settled in Beijing.
  \item \emph{Recall:} say the Mandarin for to renovate a flat, then the Mandarin for not until then, then say \emph{dìngjū} again
\end{itemize}

\begin{tcolorbox}[breakable,colback=teal!4,colframe=teal!35!black,title={Before you move on}]
What does dìngjū mean? (\textquotedblleft{}To settle down; to make a place your home\textquotedblright{}.) Say it once more.
\end{tcolorbox}
\section[shěbude]{shěbude — to hate to part with; to hate to leave}

\label{lesson:ZH-C313-shebude}



\begin{quote}
Before the new one: say the Mandarin for not until then, then the Mandarin for to make a place your home.
\end{quote}

\subsection*{You'll want to know: shěbude}
\emph{shěbude} — \textquotedblleft{}to hate to part with; to hate to leave\textquotedblright{}.

\begin{quote}
\emph{wǒ shěbude líkāi zhèr} — I hate to leave this place.
\end{quote}

Its characters come later; for now it is a word for the ear.

The tone of \zh{不} drops away here: it is said light, without a tone, shěbude.

\begin{cousinweb}
\emph{shěbude} is built from \emph{shě} 'to let go', \zh{不} 'not' and \zh{得} 'to get', all inherited from Middle Chinese, \emph{syaeX}, \emph{pjut} and \emph{tok} (Guangyun, in Baxter's transcription), which glosses the first 'to release': not able to let go.
\end{cousinweb}

\begin{grammarlens}[title={What you've built}]
One more word for telling what happened.
\end{grammarlens}

\subsection*{Your turn}
\begin{itemize}
\raggedright
  \item \emph{Say it:} \emph{shěbude}
  \item \emph{Say it:} \emph{shěbude}, once more
  \item \emph{Say it:} \emph{wǒ shěbude líkāi zhèr} — I hate to leave this place.
  \item \emph{Recall:} say the Mandarin for not until then, then the Mandarin for to make a place your home, then say \emph{shěbude} again
\end{itemize}

\begin{tcolorbox}[breakable,colback=teal!4,colframe=teal!35!black,title={Before you move on}]
What does shěbude mean? (\textquotedblleft{}To hate to part with; to hate to leave\textquotedblright{}.) Say it once more.
\end{tcolorbox}
\section[shìyìng]{shìyìng — to adapt; to adjust to new surroundings}

\label{lesson:ZH-C313-shiying}



\begin{quote}
Before the new one: say the Mandarin for to make a place your home, then the Mandarin for to hate to leave behind.
\end{quote}

\subsection*{You'll want to know: shìyìng}
\emph{shìyìng} — \textquotedblleft{}to adapt; to adjust to new surroundings\textquotedblright{}.

\begin{quote}
\emph{wǒ hái bú shìyìng zhèr de shēnghuó} — I'm not used to life here yet.
\end{quote}

Its characters come later; for now it is a word for the ear.

\begin{cousinweb}
\emph{shìyìng} is built from \emph{shì} 'to go to' and \zh{应} 'to answer', both inherited from Middle Chinese, \emph{syek} and \emph{'ingH} (Guangyun, in Baxter's transcription), which glosses the first 'to go to' and 'pleased' and the second 'to respond'.
\end{cousinweb}

\begin{grammarlens}[title={What you've built}]
One more word for telling what happened.
\end{grammarlens}

\subsection*{Your turn}
\begin{itemize}
\raggedright
  \item \emph{Say it:} \emph{shìyìng}
  \item \emph{Say it:} \emph{shìyìng}, once more
  \item \emph{Say it:} \emph{wǒ hái bú shìyìng zhèr de shēnghuó} — I'm not used to life here yet.
  \item \emph{Recall:} say the Mandarin for to make a place your home, then the Mandarin for to hate to leave behind, then say \emph{shìyìng} again
\end{itemize}

\begin{tcolorbox}[breakable,colback=teal!4,colframe=teal!35!black,title={Before you move on}]
What does shìyìng mean? (\textquotedblleft{}To adapt; to adjust to new surroundings\textquotedblright{}.) Say it once more.
\end{tcolorbox}
